\section{Feature Representation Construction and Cleaning}
\label{sec:features}

\subsection{Overview}

Using the optimal descriptor configurations from Section~\ref{sec:phase1}
(Phase~1), multiple feature representations were constructed for classifier
evaluation. This serves two purposes: to investigate the individual contribution
of each feature family, and to assess whether combining complementary descriptors
improves discrimination. The construction progressively integrates intensity
statistics, texture patterns, frequency-domain information, and structural
gradients.

\subsection{Statistical Features (A)}
\label{sec:feat_stats}

First-order statistical features summarize the global intensity distribution of
an MRI image. For each preprocessed image, eleven descriptors were extracted:
the mean, standard deviation, and variance of the intensities; the histogram
entropy (Shannon entropy over $64$ intensity bins); the skewness and kurtosis of
the intensity distribution; and five percentiles ($10^{\text{th}}$,
$25^{\text{th}}$, $50^{\text{th}}$/median, $75^{\text{th}}$, and
$90^{\text{th}}$). Together these capture the brightness level, spread, and shape
of the intensity distribution associated with different tumor types. This set is
denoted $A$.

\subsection{Texture Features (B)}
\label{sec:feat_texture}

Texture is among the most informative cues in brain tumor MRI, since tumor
tissues often exhibit distinct local patterns and structural heterogeneity. The
texture representation concatenates the optimized GLCM and LBP descriptors. The
GLCM component captures second-order spatial relationships through the six
properties selected in Phase~1 (contrast, dissimilarity, homogeneity, energy,
correlation, and ASM), while the LBP component encodes local micro-patterns as
uniform binary signatures. The fused texture set is $B = \mathrm{GLCM_{opt}}
\oplus \mathrm{LBP_{opt}}$, where $\oplus$ denotes concatenation.

\subsection{Frequency-Domain Features (C)}
\label{sec:feat_dwt}

Many tumor-related structures are discriminative across spatial frequencies. The
optimized DWT decomposes each image into sub-bands corresponding to different
frequency ranges and orientations, and statistical measures from the
approximation and detail coefficients summarize both coarse and fine structures.
This set is $C = \mathrm{DWT_{opt}}$.

\subsection{Gradient-Based Features (D)}
\label{sec:feat_hog}

Structural boundaries, edges, and shape cues help distinguish tumor classes. The
optimized HOG descriptor characterizes these properties through local
distributions of gradient orientations, focusing on the geometric and
morphological aspects of tumor regions. This set is $D = \mathrm{HOG_{opt}}$.

\subsection{Feature Fusion Strategies}
\label{sec:feat_sets}

Because individual families describe different aspects of the image, combining
them may improve discrimination through complementary information. Four fusion
representations were constructed: $A{+}B$ (statistics and texture), $B{+}C$
(texture and frequency), $A{+}B{+}C$ (statistics, texture, and frequency), and
$\mathrm{Full(opt)} = A{+}B{+}C{+}D$, which incorporates all optimized
descriptors. Table~\ref{tab:featuresets} summarizes the eight evaluated
representations together with their raw (pre-cleaning) dimensionality.

\begin{table}[t]
\centering
\caption{Evaluated feature representations. Dimensionality is reported before
cleaning; variance and correlation filtering subsequently reduce these counts.}
\label{tab:featuresets}
\renewcommand{\arraystretch}{1.2}
\begin{tabular}{@{}llc@{}}
\toprule
\textbf{Set} & \textbf{Components} & \textbf{Dim.} \\
\midrule
A            & Statistical features                         & $11$   \\
B            & Optimized GLCM + LBP                          & $20$   \\
C            & Optimized DWT                                 & $16$   \\
D            & Optimized HOG                                 & $1176$ \\
A+B          & Statistical + Texture                         & $31$   \\
B+C          & Texture + Frequency                           & $36$   \\
A+B+C        & Statistical + Texture + Frequency             & $47$   \\
Full(opt)    & Statistical + Texture + Frequency + Gradient  & $1223$ \\
\bottomrule
\end{tabular}
\end{table}

\subsection{Motivation for Feature Cleaning}
\label{sec:feat_cleaning}

Feature fusion increases dimensionality, which can introduce redundant and highly
correlated variables, noise sensitivity, additional computational cost, and a
greater risk of overfitting. A cleaning pipeline was therefore applied to each
representation prior to classifier training.

\subsection{Variance-Based Filtering}

The first stage removes constant features, whose variance
\begin{equation}
\mathrm{Var}(f) = \frac{1}{N}\sum_{i=1}^{N}(f_i - \bar{f})^2
\end{equation}
is zero across the training set ($N$ samples, feature values $f_i$, mean
$\bar{f}$). Such features carry no discriminative information and only increase
model complexity.

\subsection{Correlation-Based Reduction}

Even after variance filtering, many handcrafted descriptors remain highly
correlated and convey nearly identical information, which can harm model
stability. Pairwise Pearson correlations
\begin{equation}
\rho_{ij} = \frac{\mathrm{Cov}(f_i, f_j)}{\sigma_i \sigma_j}
\end{equation}
were computed, and for each pair with $|\rho_{ij}| > 0.95$ one feature was
removed. This pruning is applied whenever the representation is of moderate
dimensionality (below $5000$ features), which holds for all eight sets
considered here.

\subsection{Standardization}

The retained features were standardized by z-score normalization,
\begin{equation}
z = \frac{x - \mu}{\sigma},
\end{equation}
where $\mu$ and $\sigma$ are estimated on the training set only and applied
unchanged to the test set, preventing information leakage. Standardization
ensures that features on different scales contribute comparably and is
particularly important for distance- and margin-based classifiers such as KNN
and SVM.

\subsection{Final Feature Spaces}

After variance filtering, correlation reduction, and standardization, the cleaned
representations are discriminative (owing to separability-guided descriptor
selection), compact (owing to filtering), and classifier-ready (owing to
standardization). Figure~\ref{fig:cleaning} summarizes the dimensionality
reduction achieved by the cleaning pipeline; high-dimensional sets containing
HOG retain a smaller fraction of features owing to correlation pruning.
Figure~\ref{fig:tsne} visualizes class separability in the cleaned
$\mathrm{Full(opt)}$ representation via t-SNE projection. These representations
form the experimental basis for the classifier benchmarking study presented in
the next section.

\begin{figure}[t]
\centering
\begin{subfigure}[b]{0.48\linewidth}
\includegraphics[width=\linewidth]{feature_cleaning_kept_ratio}
\caption{Kept-feature ratio}
\end{subfigure}
\hfill
\begin{subfigure}[b]{0.48\linewidth}
\includegraphics[width=\linewidth]{feature_cleaning_dims}
\caption{Dimensions before/after}
\end{subfigure}
\caption{Effect of variance and correlation filtering on feature-set
dimensionality.}
\label{fig:cleaning}
\end{figure}

\begin{figure}[t]
\centering
\includegraphics[width=\linewidth]{tsne_best_feature_set}
\caption{t-SNE projection of the cleaned $\mathrm{Full(opt)}$ feature space
(training set).}
\label{fig:tsne}
\end{figure}
